\begin{helper}
Fix the actual split-outside-blocker witness surface: the graphs \(G,G'\), the
activation--priority laws \(\nu,\nu'\) with common parameter
\(\gamma/\Delta\in[0,1]\), the local embedding \(\phi\), the root \(r\), the
candidate set \(X\), the private blocker map \(\beta\), the injective outside
blocker list \(z_0,\ldots,z_{n-1}\), the candidate sets
\[
  B_k=\{x\in X:xz_k\in E(G)\},
\]
the source and target coordinate layers, the paired refined cells
\(C_\omega,C'_\omega\), the common weights \(w(\omega)\), and the complete
stage events \(S_j(\omega),T_j(\omega)\) with retained source blockers
\(i\ge j\) and already private target blockers \(i<j\).  Also fix the
adjacent boundary events \(L_k(\omega),R_k(\omega)\), and assume that they
are the displayed current common-blocker and current private-blocker boundary
events for the same \(k,\omega\) surface.  Assume the actual coordinate and
graph-replacement clauses on this same surface.  Assume also the prescribed
common refined-cell clauses: \(w(\omega)\ge0\),
\(\operatorname{ofReal}(w(\omega))=\nu(C_\omega)=\nu'(C'_\omega)\), pairwise
disjointness of the source and target cells, unit-priority support coverage
up to null complements, and the common-coordinate transport rule between
\(C_\omega\) and \(C'_\omega\).  Finally assume the displayed definitions of
\(S_j(\omega)\) and \(T_j(\omega)\) on this same cell surface.
Assume, in addition, the same-surface mixed data needed by the Claim--2
chain: a paired mixed measure surface \(\mathcal S\) of the kind defined in
\(\noderef{SamplingSplitOutsideBlockerPairedMixedMeasureSurface}\) for this
same \(G,G',\phi,r,X,n,z,B,\beta,C,C'\), and \(w\).  Assume directly the two
endpoint mass identities for \(\mathcal S\): the \(j=0\) mixed mass is the
true source survival mass in \(C_\omega\), and the \(j=n\) mixed mass is the
true target survival mass in \(C'_\omega\).  Assume also the adjacent
shared-context comparison for \(\mathcal S\): for every
\(k<n\) and \(\omega\), the two consecutive mixed masses split as
\[
M_k(\omega)=U_{k,\omega}+P^-_{k,\omega},\qquad
M_{k+1}(\omega)=U_{k,\omega}+P^+_{k,\omega},
\]
with \(P^-_{k,\omega}\le P^+_{k,\omega}\), exactly as required by
\(\noderef{SamplingSplitOutsideBlockerActualPairedSurfaceAdjacentMonotonicity}\).

Then there is a single real family \(H(j,\omega)\) such that
\[
H(j,\omega)\ge0,\qquad w(\omega)=0\Rightarrow H(j,\omega)=0,
\]
the endpoint identities are the true source and target endpoint cell masses,
\[
\operatorname{ofReal}(w(\omega)H(0,\omega))
=\nu(\{I_G(A,\pi)\cap X\ne\emptyset\}\cap C_\omega),
\]
\[
\operatorname{ofReal}(w(\omega)H(n,\omega))
=\nu'(\{I_{G'}(A',\pi')\cap \phi(X)\ne\emptyset\}\cap C'_\omega),
\]
and every adjacent replacement is monotone:
\[
  H(k,\omega)\le H(k+1,\omega)\qquad(k<n).
\]
This is the complete-stage mass law used by the hybrid producer.  It is a
one-sided aggregate comparison on the actual witness surface, not a
same-index source--target equality and not a sum of paired outside-\(B_k\)
primitive atom equalities.
\end{helper}
\begin{proof}
Let \(\mathcal S\) be the paired mixed measure surface assumed in the
statement.  By the statement of
\(\noderef{SamplingSplitOutsideBlockerPairedMixedMeasureSurface}\), this
single object contains a real family \(M_j(\omega)\) on the same
\(n,z,B,\Omega,C,C'\), and \(w\) witness surface.  It also states the two
normalization properties
\[
  M_j(\omega)\ge0,\qquad w(\omega)=0\Rightarrow M_j(\omega)=0
\]
for every \(j\) and \(\omega\), and the measured-stage identity
\[
\operatorname{ofReal}(w(\omega)M_j(\omega))
=\lambda(\widehat S_j(\omega)\cap\widehat C_\omega).
\]
Define the desired hybrid family by
\[
  H(j,\omega):=M_j(\omega).
\]
This is one global choice of \(H\), made before any endpoint or adjacent case
is checked; in particular \(H(k,\omega)\) and \(H(k+1,\omega)\) are not chosen
separately inside the \(k\)-th comparison.

The nonnegativity clause for \(H\) follows immediately from the corresponding
nonnegativity field of the paired surface:
\[
  H(j,\omega)=M_j(\omega)\ge0.
\]
Likewise, if \(w(\omega)=0\), then the zero-weight field of the same paired
surface gives \(M_j(\omega)=0\), hence \(H(j,\omega)=0\).

Next use the source endpoint mass premise assumed in the statement.  It says,
for every \(\omega\),
\[
\operatorname{ofReal}(w(\omega)M_0(\omega))
=\nu(\{(A,\pi): I_G(A,\pi)\cap X\ne\emptyset\}\cap C_\omega).
\]
Since \(H(0,\omega)=M_0(\omega)\) by definition, this is precisely the required
source endpoint identity for \(H\).  The target endpoint mass premise says,
for every \(\omega\),
\[
\operatorname{ofReal}(w(\omega)M_n(\omega))
=\nu'(\{(A',\pi'): I_{G'}(A',\pi')\cap \phi(X)\ne\emptyset\}\cap C'_\omega).
\]
Again substituting \(H(n,\omega)=M_n(\omega)\) gives the required target
endpoint identity.  These are endpoint-only uses: no statement about an
endpoint projection set or an intermediate source-only or target-only
residual projection is invoked.

Finally fix \(k<n\) and \(\omega\).  The assumed shared-context decomposition
for the same paired surface is exactly the hypothesis of
\(\noderef{SamplingSplitOutsideBlockerActualPairedSurfaceAdjacentMonotonicity}\).
Therefore that node gives
\[
  M_k(\omega)\le M_{k+1}(\omega).
\]
Replacing \(M\) by the definition of \(H\) gives
\[
  H(k,\omega)\le H(k+1,\omega).
\]
All clauses in the asserted \(H\)-family have now been verified from the one
paired surface \(\mathcal S\), its endpoint identities, and its adjacent
same-context comparisons.
\end{proof}
